\documentclass[12pt]{article}
\usepackage[utf8]{inputenc}
\usepackage[T1]{fontenc}
\usepackage{hyperref}
\title{Ukwuani-Aboh-Ndoni Language Resources for K-12 Curriculum Development}
\author{Compiled by Claude Sonnet 5 (claude-sonnet-5) for Andr\'as Kornai}
\date{}
\begin{document}
\maketitle

\section{Introduction}
Ukwuani-Aboh-Ndoni (ISO 639-3: ukw) is spoken mainly in Delta State, Nigeria, with some speakers in Rivers State. Its dialects are Ukwuani, Aboh and Ndoni. Sources describe it as Igboid. Speaker counts vary: Wikipedia cites about 550,000 from a 1973 figure, and Joshua Project counts about 367,000 people. Searches found very little material written for schools. This report lists only what the searches returned.

\section{Primary Educational Resources}
No primers, textbooks or school curricula in Ukwuani-Aboh-Ndoni turned up in the searches. Wikipedia states that the language is among the 26 Nigerian languages approved by the Federal Government for mother-tongue education. I did not find a source confirming that classroom materials exist.

\section{Linguistic and Reference Materials}
\begin{itemize}
\item Wikipedia overview of the language, with dialects and classification: \url{https://en.wikipedia.org/wiki/Ukwuani-Aboh-Ndoni_language}
\item Wikipedia on the related Onikwu community: \url{https://en.wikipedia.org/wiki/Onikwu}
\end{itemize}
No dictionary or grammar was found.

\section{Religious and Community Resources}
\begin{itemize}
\item Joshua Project people-group profile for Ukwuani-Aboh: \url{https://www.joshuaproject.net/people_groups/15699/NI}
\item Global Recordings Network listing for the language, which may hold gospel audio: \url{https://globalrecordings.net/en/language/ukw}
\end{itemize}
I did not open either page, so I cannot say what they contain beyond the listing.

\section{Digital and Audio Resources}
The only audio lead is the Global Recordings Network page above. No corpora, software or video were found.

\section{Curriculum Development Context}
Little documentation of this language turned up. A curriculum would probably start with new documentation: community orthography work, word lists and recorded speech. The searches did not identify any local organisation or school programme.
\end{document}